% TOP_PROBLEM: {"id":30007644,"problem_number":"LOCAL-30007644","title":"Drivers of extreme wealth concentration","category_id":21,"category":{"id":21,"name":"economics","display_name":"Economics","description":"Open economic theory statements, published research agendas, and proposed scoped specifications; question provenance and historical priority are qualified per record.","slug":"economics","order_index":21,"created_at":"2026-10-08T00:00:00Z"},"status":"research_question","external_url":null,"legacy_ids":[],"published":false,"statement":"In the explicitly specified setting: What relative roles do unequal investment returns, saving behavior, entrepreneurial risk, and inherited resources play in generating the extreme upper tail of wealth? The mathematical statement above fixes the target, assumptions and scope of this catalogue formulation.","economics":{"original_id":133,"field":"Inequality and economic mobility","kind":"identification","formalization_basis":"adapted_research_question","source_status":"research_agenda","priority_status":"earliest_found","priority_note":"Earliest directly inspected qualifying passage in this audit; earlier origins were not established. A review or research-agenda statement does not imply original priority.","earliest_verified_reference":{"authors":["Mariacristina De Nardi","Giulio Fella"],"title":"Saving and Wealth Inequality","year":2017,"url":"https://gabriel-zucman.eu/files/teaching/DeNardiFella17.pdf","doi":"10.1016/j.red.2017.06.002","locator":"CEPR DP 11746; introduction, printed pp. 2–4; section 8.1, printed pp. 26–27","evidence_summary":"The review calls for richer data to distinguish wealth mechanisms and further understanding of endogenous heterogeneous returns."},"current_reference":{"authors":["Mariacristina De Nardi","Giulio Fella"],"title":"Saving and Wealth Inequality","year":2017,"url":"https://gabriel-zucman.eu/files/teaching/DeNardiFella17.pdf","doi":"10.1016/j.red.2017.06.002","locator":"CEPR DP 11746; introduction, printed pp. 2–4; section 8.1, printed pp. 26–27","evidence_summary":"The review calls for richer data to distinguish wealth mechanisms and further understanding of endogenous heterogeneous returns."},"formalization_note":"The cited passage establishes a research direction, not this exact mathematical statement. The notation, population, policy menu, assumptions, and estimands are an original adaptation for this catalogue; no universal unsolved theorem or source priority is claimed."},"statusQualification":"Published question or research agenda verified at the cited locator. The mathematical specification is adapted for this catalogue and includes additional assumptions; source publication does not establish first-ever priority or certify this exact model remains unsolved. The cited passage establishes a research direction, not this exact mathematical statement. The notation, population, policy menu, assumptions, and estimands are an original adaptation for this catalogue; no universal unsolved theorem or source priority is claimed.","statusReviewedAt":"2026-10-08"}
% FIRST_POSTED: 2026-10-08
% ENTRY_KIND: statement_only
% !TeX program = lualatex
\documentclass[11pt]{article}
\usepackage{fontspec}
\usepackage[a4paper,margin=25mm]{geometry}
\usepackage{amsmath,amssymb,mathtools}
\usepackage{xurl}
\usepackage[hidelinks]{hyperref}
\newcommand{\sref}[2]{\hyperlink{source:#1:#2}{[S#2]}}
\setlength{\emergencystretch}{3em}
\begin{document}
% problemId:30007644
\section*{Drivers of extreme wealth concentration}\label{30007644}
\subsection{Definitions and mathematical statement}
Consider a specified economy of households whose wealth evolves according to \(W_{i,t+1}=(1+r_{it})W_{it}+Y_{it}-C_{it}+B_{it}\), where \(r_{it}\) is the realized portfolio return, \(Y_{it}\) labor and entrepreneurial income, \(C_{it}\) consumption, and \(B_{it}\) net transfers including inheritances. Let \(\mathcal M\) be a declared class of models governing saving, portfolio choice, entrepreneurship, and intergenerational transfers. For a top population share \(p\in(0,1)\), assuming positive finite total wealth under every compared model, let \(S_{p,t}(M)\) be the model-implied share of total wealth owned by that group. Identify which models jointly reproduce concentration, individual wealth transitions, returns, portfolios, and inheritance histories, rather than merely matching \(S_{p,t}\). For each mechanism \(k\), define a counterfactual \(M^{-k}\) that removes that mechanism according to a stated intervention while recomputing equilibrium. Estimate \(\Delta_{k,t}=S_{p,t}(M)-S_{p,t}(M^{-k})\). Since mechanisms interact, report the counterfactual ordering or a specified averaging rule and avoid interpreting contributions as an invariant additive identity. The unresolved quantitative task is to distinguish relative importance and validate predictions outside fitted moments. Assumptions about unobserved skill, bequest motives, and endogenous returns must be exposed and tested where possible.

\par\textbf{Formulation and scope.} The cited passage establishes a research direction, not this exact mathematical statement. The notation, population, policy menu, assumptions, and estimands are an original adaptation for this catalogue; no universal unsolved theorem or source priority is claimed.

\par\textbf{Open-question provenance.} An explicit question or research agenda was verified at the source locator. This does not by itself certify that every later version remains unresolved \sref{30007644}{1}.

\par\textbf{Historical priority.} Earliest directly inspected qualifying passage in this audit; earlier origins were not established. A review or research-agenda statement does not imply original priority.

\subsection{Short English statement}
In the explicitly specified setting: What relative roles do unequal investment returns, saving behavior, entrepreneurial risk, and inherited resources play in generating the extreme upper tail of wealth? The mathematical statement above fixes the target, assumptions and scope of this catalogue formulation.

\subsection{Sources}
\begin{itemize}\raggedright
\item[\textnormal{[S1]}]\hypertarget{source:30007644:1}{} Mariacristina De Nardi, Giulio Fella. 2017. Saving and Wealth Inequality. CEPR DP 11746; introduction, printed pp. 2–4; section 8.1, printed pp. 26–27.
\url{https://gabriel-zucman.eu/files/teaching/DeNardiFella17.pdf}.
DOI: 10.1016/j.red.2017.06.002.
{\small\textit{Earliest verified question/agenda reference; historical priority qualified below. The review calls for richer data to distinguish wealth mechanisms and further understanding of endogenous heterogeneous returns.}}
\end{itemize}
\end{document}
